\section{Problem Statement}
% In modern software development, data is at the core of nearly every application. However, as database technologies continue to diversify—particularly with the coexistence of SQL and NoSQL systems—developers are increasingly burdened with managing multiple tools, interfaces, and workflows. This fragmentation creates significant challenges in efficiency, learning, collaboration, and scalability.

% Traditional database management requires switching between different platforms for tasks such as schema design, query execution, data manipulation, and visualization. SQL databases (e.g., PostgreSQL, MySQL) and NoSQL models (e.g., MongoDB) typically rely on entirely separate tools with unique configurations, syntaxes, and interfaces. As a result, developers, especially students and early-stage engineers, must invest substantial time and effort just to understand, configure, and operate these systems. This leads to steep learning curves, frequent errors, inconsistent workflows, and slower project development cycles.

% Modern development environments increasingly rely on AI to automate repetitive tasks, assist with query formulation, interpret complex schemas, and accelerate onboarding for developers. However, most existing DBMS tools do not integrate AI capabilities in a meaningful way. Developers dealing with unfamiliar schemas or large datasets must spend excessive time understanding relational structures, dependencies, and data flow manually. Query writing—especially for complex joins, aggregations, or NoSQL pipeline queries—becomes time-consuming and error-prone.

% Additionally, while cloud-based database offerings exist, they typically focus on hosting rather than providing an integrated development workspace. They do not combine schema design, real-time editing, AI-driven query support, version control, and multi-database management into a single cohesive environment. This gap creates inefficiencies for teams that want a reliable, centralized, AI-powered platform for building and managing database projects.

% Finally, the demand for intelligent, natural-language interfaces is rising rapidly. Users expect systems that can understand plain language requests such as "generate a schema for an e-commerce store", "show customers who purchased more than five items", or "explain the relationship between these tables". Yet existing platforms rarely leverage Retrieval-Augmented Generation (RAG) to provide context-aware assistance tailored to the specific database structure and stored data.

% Therefore, there is a clear need for a unified, cloud-based Database-as-a-Service (DBaaS) platform that integrates SQL and NoSQL database management, provides intuitive schema design tools, supports AI-assisted interactions, and eliminates the fragmented, inefficient workflows developers face today. The absence of such a system results in slower development cycles, higher cognitive load, reduced accessibility for beginners, and limited adoption of best practices in database engineering.

Modern software development relies heavily on data, yet the diversity of database technologies, particularly the coexistence of SQL and NoSQL systems, imposes significant challenges on developers. Managing multiple tools and interfaces for schema design, query execution, data manipulation, and visualization creates fragmented workflows, steep learning curves, and increased error rates.

While AI can automate repetitive tasks, assist with query formulation, and accelerate onboarding, most existing database tools lack meaningful AI integration. Similarly, cloud-based offerings focus on hosting rather than providing a unified development environment.

This underscores the need for a cloud-based, AI-powered Database-as-a-Service (DBaaS) platform that integrates SQL and NoSQL management, intuitive schema design, and intelligent query support, streamlining workflows and improving accessibility for both novice and experienced developers.